\chapter{Implémentation}

\section{Organisation générale du prototype}

Le prototype \logminer{} est construit autour d'un noyau développé en Python et d'un tableau de bord web destiné à la consultation des résultats. Les principaux composants applicatifs sont regroupés dans le paquet \texttt{src/logminer}, tandis que les scripts utilisés pour les expérimentations sont placés dans le dossier \texttt{scripts}. Les résultats consolidés sont enregistrés dans \texttt{data/processed} avant d'être utilisés pour produire les tableaux, les figures et les captures présentés dans ce mémoire.

Cette organisation permet de séparer le code applicatif, les expérimentations et les artefacts utilisés pour la rédaction. Elle facilite également la vérification des résultats, puisque les valeurs présentées dans le mémoire peuvent être reliées aux scripts et aux fichiers qui les ont produites.

\section{Outils et environnement}

L'implémentation repose principalement sur Python pour le développement du pipeline, scikit-learn pour les modèles d'apprentissage automatique légers et FastAPI pour l'exposition des services locaux. Une interface web complète cet ensemble afin de permettre la consultation et l'interaction avec les résultats.

Les données intermédiaires et les résultats sont principalement conservés sous forme de fichiers CSV et JSONL. Les sorties destinées au mémoire sont ensuite transformées en tableaux, figures et captures.

Cette pile technologique a été retenue pour trois raisons principales : sa facilité de mise en œuvre, sa compatibilité avec un environnement local et sa simplicité de reproduction. Elle permet également de modifier rapidement une partie du prototype sans nécessiter une infrastructure logicielle particulièrement lourde.

\section{Pipeline de traitement}

Le pipeline constitue le point central de l'implémentation. Il organise les différentes étapes nécessaires pour passer d'une source de journaux brute à un ensemble d'informations pouvant être analysées.

Les principales étapes sont les suivantes :

\begin{enumerate}

\item lecture ou découverte des sources de journaux ;

\item parsing des formats reconnus ou conservation des entrées lorsque leur format ne peut pas être identifié ;

\item normalisation des informations extraites dans un schéma commun ;

\item routage vers une famille de journaux ;

\item détection ou classification à l'aide du modèle correspondant ;

\item corrélation des anomalies candidates ;

\item export des résultats au format CSV et journalisation des informations nécessaires à l'audit.

\end{enumerate}

Le module \texttt{pipeline.py} coordonne principalement cette chaîne de traitement dans sa version locale. Le fichier \texttt{api.py} expose les principales fonctions au moyen de FastAPI. Les modules \texttt{model\_router.py} et \texttt{model\_compare.py} sont utilisés respectivement pour le routage des données vers les familles appropriées et pour la comparaison de certains modèles.

Dans cette implémentation, le pipeline ne joue pas uniquement le rôle d'un simple enchaînement de fonctions. Il constitue également une couche d'intégration entre le code, les expérimentations et les résultats utilisés dans le mémoire.

Les fichiers produits doivent ainsi conserver une structure suffisamment stable pour être réutilisés par les scripts d'évaluation. Les colonnes attendues doivent rester cohérentes, les chemins doivent pouvoir être retrouvés et les erreurs doivent fournir des informations suffisamment explicites pour être diagnostiquées.

Ce besoin explique notamment le recours fréquent aux formats CSV et JSONL. Ils sont moins spécialisés qu'une base de données dédiée, mais présentent l'avantage d'être faciles à lire, à comparer et à archiver. Ils conviennent donc bien à un prototype dont les résultats doivent pouvoir être examinés manuellement.

Une séparation est également maintenue entre les données brutes, les données transformées et les artefacts utilisés dans la rédaction.

Les fichiers bruts volumineux peuvent être conservés en dehors du système de versionnement. Les résultats utiles aux expérimentations sont ensuite exportés vers les dossiers prévus à cet effet. Les tableaux et les figures utilisés dans le mémoire sont finalement produits à partir de ces sorties.

Cette organisation permet de réduire le risque de présenter une valeur numérique qui ne serait plus directement reliée à l'expérience dont elle provient.

\section{Parsing multiformat}

Le prototype doit pouvoir traiter plusieurs familles de données qui n'utilisent pas les mêmes formats.

Les formats pris en charge comprennent notamment CSV, JSONL, syslog, Apache, CloudTrail, HDFS, BGL, Windows XML ou EVTX exporté, CEF/LEEF, tcpdump sous forme textuelle ainsi que certains fichiers pcap lorsque les dépendances nécessaires sont disponibles.

La stratégie retenue pour le parsing est volontairement conservatrice.

Lorsqu'un format est correctement reconnu, le parseur extrait les informations qui peuvent être structurées. Lorsque la reconnaissance n'est que partielle, les champs identifiés sont conservés tandis que le message original reste disponible.

Lorsqu'une ligne est incomplète, inconnue ou partiellement corrompue, elle n'est pas automatiquement supprimée.

Cette décision est importante dans une chaîne d'analyse de sécurité. Ignorer silencieusement un événement peut créer un angle mort. À l'inverse, conserver toutes les données sans distinction pourrait produire une quantité importante de bruit.

Le compromis adopté consiste donc à structurer ce qui peut être interprété tout en conservant les informations restantes sous forme de message brut ou de métadonnées.

\begin{table}[H]
\centering
\caption{Exemples de formats pris en charge et traitement attendu}
\label{tab:formats-implementation}
\scriptsize

\begin{tabularx}{\textwidth}{p{3cm}p{3cm}X}
\toprule
\textbf{Format} & \textbf{Type de source} & \textbf{Traitement dans le prototype} \\
\midrule

Windows XML/EVTX exporté &
Système local &
Extraction des événements, des canaux, des messages et des champs de sécurité disponibles. \\

Syslog/Linux auth &
Système Unix/Linux &
Extraction de l'horodatage, de l'hôte, du processus, de l'utilisateur et du message d'authentification lorsqu'ils sont présents. \\

Apache &
Application web &
Extraction de la méthode, du chemin, du code HTTP, de l'adresse du client et du volume lorsque ces informations peuvent être récupérées. \\

Wazuh export &
HIDS/SIEM &
Conservation des informations enrichies telles que la règle, l'agent, le groupe, le niveau de sévérité et le message. \\

HDFS/BGL &
Journaux systèmes publics &
Normalisation des messages et préparation des données pour les expérimentations. \\

CICIDS/CIC-DDoS2019 &
Flux réseau tabulaire &
Exploitation des caractéristiques numériques et des labels disponibles dans les expérimentations supervisées. \\

Format inconnu &
Cas limite &
Conservation de l'entrée dans une route de fallback avec le message original et un statut explicite. \\

\bottomrule
\end{tabularx}
\end{table}

\section{Registre des modèles}

Les modèles entraînés sont enregistrés dans le dossier \texttt{models}. Un registre permet d'associer chaque modèle à la famille de journaux concernée, au type d'apprentissage utilisé et aux principales informations nécessaires à son exploitation.

Cette organisation répond à une difficulté fréquente dans les prototypes d'apprentissage automatique : le terme ``modèle'' peut désigner l'algorithme, le fichier entraîné ou encore une configuration expérimentale particulière.

Dans \logminer{}, un modèle utilisable doit être associé à une famille de journaux, à un type d'apprentissage, à un artefact sauvegardé, aux caractéristiques qu'il attend et, lorsque celles-ci sont disponibles, aux métriques obtenues lors de son évaluation.

Cette information facilite le routage et réduit le risque d'utiliser un modèle sur des données incompatibles.

Les modèles sont notamment sauvegardés sous forme d'artefacts \texttt{joblib}. Cette solution évite de relancer l'entraînement lors de chaque démonstration ou campagne de test.

Elle permet donc de stabiliser les expérimentations et d'en réduire le coût d'exécution.

Cette facilité impose toutefois une discipline de documentation. Un artefact doit rester associé à son jeu de données d'origine, à sa date de génération, au type de caractéristiques utilisées, à ses résultats de validation et à ses principales limites.

Sans ces informations, un fichier de modèle peut rester techniquement utilisable tout en devenant difficile à interpréter scientifiquement.

\section{Agents intelligents multitâches}

La version avancée du prototype introduit un environnement d'exécution dédié aux agents dans le fichier \texttt{agents/intelligent\_runtime.py}.

Chaque agent déclare les tâches qu'il peut exécuter et conserve un état interne décrivant notamment ses succès, ses erreurs, certaines durées d'exécution et ses décisions récentes.

Un mécanisme de heartbeat permet également de suivre son état.

Le choix des tâches repose sur une politique explicite combinant plusieurs éléments : priorité, confiance, coût, fiabilité observée, urgence, nombre de tentatives précédentes et diversité des tâches récemment exécutées.

Les tâches concernées couvrent notamment la découverte des sources, le parsing, le routage des modèles, la détection et la corrélation.

Deux modes d'exécution sont disponibles.

Le premier utilise une file locale en mémoire. Ce mode est principalement adapté aux démonstrations et aux tests simples.

Le second repose sur Redis Streams et ses consumer groups. Plusieurs processus peuvent alors consommer un même flux de tâches.

Le worker \texttt{scripts/logminer\_intelligent\_agent\_worker.py} prend également en charge la récupération de certaines tâches non acquittées à l'aide de \texttt{XAUTOCLAIM}. Ce mécanisme permet notamment de tester le comportement du système lorsqu'un worker s'interrompt avant d'avoir confirmé le traitement d'une tâche.

Cette implémentation donne une traduction concrète aux notions d'agents, de fonctionnement multitâche et de distribution utilisées dans le mémoire.

Les agents ne sont plus uniquement des fonctions exécutées dans une séquence fixe. Ils disposent d'un ensemble de capacités déclarées, d'un état observable, d'une file de tâches et d'un mécanisme de reprise.

La distribution évaluée reste toutefois locale. Les expériences utilisent plusieurs processus indépendants communiquant à travers un bus Redis persistant. Elles ne constituent donc pas une validation complète d'un déploiement industriel réparti entre plusieurs infrastructures distantes.

\section{Mémoire adaptative implémentée}

La mémoire adaptative est introduite progressivement à plusieurs niveaux du prototype.

Dans \texttt{agents/intelligent\_runtime.py}, la structure \texttt{AgentMemory} conserve notamment les succès et les erreurs associés aux différents types de tâches, les durées récentes, les derniers messages d'erreur, les tâches exécutées et les raisons enregistrées lors de certaines décisions.

Ces informations sont sauvegardées dans des fichiers JSON afin de rester lisibles et de pouvoir être réutilisées au cours de plusieurs cycles d'exécution.

Le superviseur possède également sa propre mémoire persistante dans \texttt{agents/supervisor\_agent.py}.

La structure \texttt{SupervisorMemory} conserve les chemins déjà traités, les familles rencontrées récemment, le nombre de cycles exécutés, les succès, les erreurs et la dernière décision prise.

Lorsqu'il doit sélectionner une nouvelle source, le superviseur peut ainsi pénaliser les chemins traités récemment et conserver les raisons ayant conduit à sa sélection.

Cette logique constitue une première forme d'adaptation comportementale. Le système ne répète pas systématiquement la même décision lorsqu'un historique pertinent est disponible.

Un troisième niveau concerne les décisions prises par l'analyste.

L'endpoint \texttt{/alerts/decision} permet d'enregistrer plusieurs actions sur une anomalie candidate. Celle-ci peut notamment être acceptée, rejetée ou reclassée.

L'action, la cible, la catégorie, la sévérité et la justification éventuelle sont ajoutées au journal d'audit.

Au début d'un nouveau cycle, le superviseur peut relire ces informations et consolider certains rejets ou reclassements selon leur catégorie et leur sévérité. Ces données sont ensuite utilisées comme éléments contextuels lors de certaines décisions.

Les retours de l'analyste remplissent donc plusieurs fonctions. Ils améliorent la traçabilité, peuvent être consultés dans le tableau de bord et forment progressivement une mémoire pouvant servir à identifier les anomalies fréquemment rejetées.

Cette mémoire pourrait également préparer un futur mécanisme d'apprentissage actif hors ligne. Dans l'implémentation actuelle, elle ne déclenche cependant pas automatiquement le réentraînement des modèles.

\section{API FastAPI}

La deuxième version du prototype expose plusieurs fonctions au moyen d'une API FastAPI locale.

Cette API permet notamment de lancer certaines analyses, de consulter l'état des services, d'interroger le registre des modèles, de récupérer des informations sur les ressources et de consulter les événements enregistrés dans l'audit.

Le tableau de bord utilise également cette API pour récupérer les données nécessaires à son fonctionnement.

Cette architecture introduit une séparation entre l'interface graphique et le moteur d'analyse.

Le tableau de bord n'a donc pas besoin de connaître directement l'organisation interne des fichiers ou des modules Python.

Cette séparation rend l'architecture plus flexible. Elle facilite notamment l'ajout futur de nouveaux workers, de files d'attente, de services intermédiaires ou de connecteurs vers d'autres plateformes.

\section{Tableau de bord Ariel Logminer}

Le tableau de bord Ariel \logminer{} constitue l'interface principale de consultation et d'interaction avec le prototype.

Il regroupe les indicateurs globaux, les alertes considérées comme prioritaires, les anomalies candidates, les incidents, plusieurs graphiques de suivi, les informations relatives aux ressources, l'état des services de la V2 ainsi que le journal d'audit.

\begin{figure}[H]
\centering
\includegraphics[width=\textwidth]{dashboard_vue_ensemble.png}
\caption{Vue d'ensemble du tableau de bord Ariel \logminer{}.}
\label{fig:dashboard-overview}
\end{figure}

Une vue spécifique est consacrée aux résultats. Elle permet notamment d'examiner les anomalies candidates, d'accéder aux détails d'un incident et de consulter ou d'enregistrer les décisions de l'analyste.

\begin{figure}[H]
\centering
\includegraphics[width=\textwidth]{dashboard_resultats_detail_incident.png}
\caption{Vue des résultats, détail d'un incident et décisions de l'analyste.}
\label{fig:dashboard-results}
\end{figure}

Une autre vue est davantage orientée vers le fonctionnement technique du prototype. Elle présente notamment les ressources, les services et plusieurs informations issues de l'audit.

\begin{figure}[H]
\centering
\includegraphics[width=\textwidth]{dashboard_ressources_audit.png}
\caption{Vue technique : ressources, services et audit.}
\label{fig:dashboard-technical}
\end{figure}

\section{Audit et décisions de l'analyste}

Le tableau de bord ne se limite pas à la consultation des anomalies candidates.

Il permet également à l'analyste d'enregistrer une décision. Une anomalie peut être validée, rejetée ou reclassée.

Ces actions sont conservées dans le système d'audit afin de pouvoir être consultées ultérieurement.

Cette fonctionnalité prépare une boucle d'amélioration dans laquelle les retours humains pourraient progressivement contribuer à réduire certains faux positifs.

La boucle reste toutefois volontairement contrôlée.

Une décision prise dans le tableau de bord ne modifie pas automatiquement les paramètres d'un modèle déjà entraîné.

Elle enrichit plutôt une trace persistante qui peut ensuite être analysée afin d'identifier les catégories régulièrement rejetées, les familles produisant trop de bruit ou les seuils qui pourraient nécessiter un recalibrage.

Cette séparation permet d'améliorer progressivement le système tout en maintenant une trace des décisions ayant conduit aux adaptations envisagées.

\section{Problèmes rencontrés}

Plusieurs difficultés ont été rencontrées au cours de l'implémentation.

La première est directement liée à l'hétérogénéité des données. Certaines sources utilisent des structures tabulaires, d'autres sont essentiellement textuelles, tandis que d'autres encore contiennent déjà des informations enrichies par une plateforme de sécurité.

La deuxième difficulté concerne les modèles non supervisés. Lorsqu'aucun label fiable n'est disponible, un score élevé d'anomalie ne peut pas être présenté comme la preuve d'une attaque. Il reste un indicateur qui nécessite une interprétation supplémentaire.

La troisième difficulté concerne le tableau de bord.

Une interface trop synthétique ne donne pas suffisamment de contexte à l'analyste. Une interface qui affiche directement toutes les données techniques peut, au contraire, devenir difficile à exploiter.

Le développement a donc nécessité un compromis entre richesse des informations et lisibilité.

\section{Reproductibilité}

La reproductibilité constitue un élément important de l'implémentation.

Les principales campagnes expérimentales sont accompagnées de scripts et de fichiers de résultats permettant de retrouver les mesures utilisées dans le mémoire.

Parmi les scripts principaux figurent :

\begin{itemize}

\item \texttt{benchmark\_realtime\_workflow.py} ;

\item \texttt{run\_resource\_campaign.py} ;

\item \texttt{run\_family\_routing\_ablation.py} ;

\item \texttt{generate\_memoire\_assets.py} ;

\item \texttt{export\_memoire\_figures.py}.

\end{itemize}

Ils permettent notamment de reproduire plusieurs mesures, de consolider les résultats et de générer une partie des figures et tableaux utilisés dans le document.

\section{Organisation des modules}

Le paquet \texttt{src/logminer} contient les composants principaux du prototype.

Le pipeline coordonne les traitements. Le parseur identifie les formats et construit les représentations intermédiaires. Le routeur détermine les familles de journaux. Le détecteur applique les modèles disponibles. Le corrélateur rapproche plusieurs anomalies candidates. Enfin, le système d'audit conserve les informations liées aux exécutions et aux décisions.

Cette organisation traduit directement l'architecture modulaire décrite dans le chapitre précédent.

Le fichier \texttt{model\_compare.py} est utilisé dans l'environnement expérimental pour comparer certains modèles et consolider leurs résultats.

Il s'intègre dans une logique de registre où chaque détecteur reste associé à une famille de données, à un type d'apprentissage et à un ensemble de métriques.

Cette discipline limite les risques de comparer ou d'utiliser des résultats provenant de contextes expérimentaux différents sans le signaler.

\section{Préparation des données}

La préparation des données constitue une étape indispensable avant l'application des modèles.

Les traitements dépendent fortement du type de source.

Pour les jeux de données tabulaires, il faut notamment vérifier les colonnes disponibles, identifier le label, traiter les variables lorsque cela est nécessaire et construire les ensembles d'entraînement et de test selon le protocole expérimental retenu.

Pour les journaux textuels, le traitement porte plutôt sur l'extraction des horodatages, des messages, des sources, des utilisateurs ou encore des niveaux de sévérité lorsqu'ils sont disponibles.

Ces transformations ne sont pas considérées comme une étape invisible de l'expérimentation.

Les scripts conservés avec le projet documentent les principales opérations réalisées. Cette organisation permet de relier les tableaux, les figures et les valeurs numériques aux fichiers qui ont servi à les produire.

\section{Implémentation du tableau de bord}

Le tableau de bord a été conçu comme une véritable interface de travail et non comme une simple page de présentation.

Il regroupe plusieurs catégories d'informations : indicateurs généraux, résultats, détails des incidents, ressources disponibles, services et événements d'audit.

Plusieurs captures sont intégrées au mémoire afin de représenter ces différents usages.

La vue d'ensemble présente l'état général du prototype.

La vue consacrée aux résultats permet de parcourir les anomalies candidates et de consulter les décisions prises par l'analyste.

La vue technique présente les ressources disponibles, l'état des services et plusieurs informations issues de l'audit.

L'intérêt principal de cette interface est de rendre les sorties du pipeline plus faciles à interpréter.

Un fichier CSV peut être suffisant pour vérifier automatiquement un résultat ou réaliser une analyse statistique. Il est cependant moins adapté lorsqu'un utilisateur doit examiner plusieurs événements, comprendre leur contexte et prendre une décision.

Le tableau de bord constitue donc une couche d'interprétation entre la sortie technique du pipeline et l'analyse humaine.

\section{Choix de légèreté}

Le prototype privilégie volontairement des technologies relativement légères et accessibles.

Python et FastAPI permettent de développer rapidement les composants nécessaires. Les fichiers CSV et JSONL simplifient l'export et l'audit. Les modèles issus de scikit-learn demandent généralement moins de ressources que des architectures profondes complexes et restent relativement simples à entraîner et à déployer dans un environnement local.

Ce choix ne signifie pas que des méthodes plus avancées seraient inutiles.

Il traduit plutôt le périmètre du mémoire : valider d'abord une chaîne complète, fonctionnelle et reproductible avant d'introduire des modèles dont la complexité augmenterait également les exigences de calcul, de données et d'évaluation.

\section{Validation fonctionnelle}

La validation fonctionnelle complète les métriques quantitatives présentées dans les chapitres consacrés aux résultats.

Elle repose sur plusieurs vérifications : exécution correcte des scripts expérimentaux, production des fichiers CSV attendus, génération des figures, fonctionnement du tableau de bord, cohérence des captures et correspondance entre les tableaux générés et les résultats expérimentaux.

Ces vérifications montrent que le prototype ne se limite pas à produire des scores isolés.

Il génère également des artefacts pouvant être consultés, archivés et utilisés pour l'analyse technique ou pour la rédaction scientifique.

\section{Gestion des dépendances et portabilité}

La portabilité du prototype dépend en partie de la manière dont ses dépendances sont organisées.

Les bibliothèques indispensables au fonctionnement général sont distinguées des dépendances plus lourdes ou spécifiques à certains formats.

Cette séparation permet d'installer rapidement le pipeline principal sans imposer immédiatement toutes les bibliothèques nécessaires aux traitements spécialisés.

La lecture directe de fichiers EVTX peut, par exemple, nécessiter une bibliothèque spécifique, alors que des événements Windows exportés en XML peuvent être traités avec des outils plus standards.

Cette distinction permet de maintenir un noyau relativement léger tout en conservant la possibilité d'ajouter certaines fonctionnalités lorsque l'environnement l'exige.

La portabilité concerne également les systèmes d'exploitation.

Le projet a été principalement développé et testé dans un environnement Windows. Une partie importante des composants repose cependant sur Python et reste donc portable vers d'autres plateformes.

Les scripts PowerShell sont principalement utilisés pour certaines opérations de collecte propres à Windows, tandis que les scripts Python prennent en charge l'essentiel du traitement, de l'évaluation et de la génération des artefacts.

Cette séparation facilite une adaptation progressive du prototype vers un environnement Linux ou vers un serveur dédié.
